\chapter{What Cryptography Is}

As previewed in the preface, we start with the central problem: how do you share a secret when anyone might be watching? Imagine you write a note, fold it, and slide it across a crowded room to a friend. Between your hand and theirs, a dozen people could read it, copy it, change a word, or replace the note entirely with one of their own. Cryptography is the science of making that slip of paper behave as if the room were empty -- as if only your friend could read it, any tampering would be immediately detected, and a forged note could be told apart from a real one. It does this not with locks and walls but with mathematics, and the foundational claim it makes is that a few hundred bits of carefully chosen numbers can defeat an adversary with more computing power than exists on Earth.

This chapter builds the vocabulary and the worldview you will use for the rest of the book. We will not yet encrypt anything seriously. Instead we pin down four goals that cryptography chases (confidentiality, integrity, authentication, non-repudiation), meet the cast of characters who show up in every protocol story, learn the one principle that separates real cryptography from wishful thinking, and -- most importantly -- learn what the word ``secure'' is even allowed to mean. That last point is subtle enough that whole schemes which ``felt'' secure to clever people have collapsed embarrassingly. By the end you should have the security mindset: a reflex to ask, before trusting anything, ``secure against whom, given what, and proven how?''

\chapterline

\section{Four Goals: What Are We Actually Protecting?}
\label{ch01:four-goals}

When people first hear ``cryptography'' they think ``secret codes'' -- hiding the contents of a message. That is one goal, but only one, and historically not even the one that gets people killed most often. A message that nobody can read is useless if an attacker can swap it for a different message, or impersonate the sender, or convince a court they never sent something they did. Cryptography pursues four distinct goals, and a real system usually needs several of them at once.

\begin{description}
  \item[Confidentiality.] Only the intended recipient learns the contents. An eavesdropper who sees every bit on the wire learns nothing useful about the message (ideally, nothing at all beyond its length). This is the ``secret writing'' goal, the one encryption provides.
  \item[Integrity.] The recipient can detect if the message was altered in transit -- even by one bit, even by an attacker who controls the entire network. Note this does not mean the attacker cannot \emph{change} the bits; on an open network they always can. It means the recipient will \emph{notice} and reject the change.
  \item[Authentication.] The recipient can verify who the message came from. ``This really is from Alice, not from someone pretending to be Alice.'' There are two flavors: \emph{entity authentication} (you are talking to the right party right now) and \emph{message} or \emph{data-origin authentication} (this particular message originated with that party).
  \item[Non-repudiation.] The sender cannot later deny having sent the message. If Alice signs a contract, she cannot turn around and tell a judge ``that wasn't me.'' This is a stronger and stranger property than the others, because it must convince a \emph{third party} -- not just the recipient.
\end{description}

\begin{intuitionbox}
A useful way to keep these straight is to ask what kind of lie each one prevents.
\emph{Confidentiality} stops the lie ``I read your message.'' \emph{Integrity} stops the lie ``the message said X'' when it actually said Y. \emph{Authentication} stops the lie ``I am Alice.'' \emph{Non-repudiation} stops the lie ``I never said that'' -- told by Alice herself, after the fact. The first three protect the communicating pair against an outsider. The fourth protects each member of the pair against the \emph{other}, and against future disputes.
\end{intuitionbox}

\begin{whybox}
\textbf{Why non-repudiation is the odd one out.} The first three goals can be achieved with a shared secret key: if only Alice and Bob know $k$, then a valid authentication tag on a message proves \emph{to Bob} that Alice sent it -- nobody else could have made the tag. But that same proof is worthless \emph{to a judge}, because Bob knows $k$ too, so Bob could have forged the message himself. With shared-key authentication, Alice can always plausibly say ``Bob faked it,'' and the judge cannot tell who is lying. Non-repudiation therefore demands an \emph{asymmetry} the other goals do not: Alice must hold a secret that Bob does \emph{not} have, so that a valid tag could only have come from Alice. That is exactly what a digital signature provides -- a private signing key only Alice holds, and a public verification key everyone (including the judge) can check. This is the first hint that some goals are unreachable with symmetric cryptography alone, foreshadowing \S\ref{ch01:sym-vs-asym}.
\end{whybox}

These goals are independent. You can have any one without the others, and confusing them is a classic source of broken systems. A WiFi network that encrypts your traffic (confidentiality) but lets an attacker flip bits undetected (no integrity) can be manipulated in nasty ways -- and historically was. An email that is signed (authentication, non-repudiation) but sent in the clear (no confidentiality) is perfectly readable by anyone on the path. Encryption alone gives you neither integrity nor authentication; that surprises newcomers, and we will hammer on it repeatedly.

\begin{trapbox}
\textbf{``I encrypted it, so it can't be tampered with.''} False, and dangerously so. Most basic encryption schemes are \emph{malleable}: an attacker who cannot read the plaintext can still flip chosen bits of the ciphertext and cause predictable changes to the decrypted plaintext. We will see in Chapters~9 and~10 that flipping bit $i$ of a one-time-pad ciphertext flips bit $i$ of the plaintext exactly. Confidentiality and integrity are separate goals requiring separate mechanisms. Modern designs bundle them with \emph{authenticated encryption}, which we treat later -- but the lesson starts here.
\end{trapbox}

\begin{applicationbox}
When you load a banking site over HTTPS, all four goals are in play. \emph{Confidentiality}: nobody on the coffee-shop WiFi reads your password. \emph{Integrity}: the network cannot alter the page to inject a fake login form without detection. \emph{Authentication}: the certificate proves you are talking to the real bank, not an impostor. \emph{Non-repudiation} shows up less in the browser and more in the back end -- digitally signed transaction logs that the bank can present if you dispute a transfer. TLS, the protocol behind HTTPS, is essentially an engineering answer to ``how do we get the first three over a hostile network, starting from nothing but a trusted list of certificate authorities?'' We will build up to it.
\end{applicationbox}

\section{What Cryptography Is Not}
\label{ch01:what-its-not}

It sharpens the definition to fence off three neighbors that people routinely confuse with cryptography.

\textbf{Cryptography is not steganography.} Steganography \emph{hides the existence} of a message -- invisible ink, data tucked in the low bits of an image, a microdot. Cryptography hides the \emph{contents} while making no secret that a message exists; an eavesdropper plainly sees that ciphertext is flying by, just not what it says. The two are complementary and orthogonal. Steganography fails the moment the hiding place is discovered (and per Kerckhoffs we must assume it will be); cryptography is designed to stay secure even when the adversary sees every encrypted bit. Serious systems encrypt first and, only if hiding the existence of communication also matters, steganographically conceal the ciphertext as a second, weaker layer.

\textbf{Cryptography is not encoding.} Base64, ASCII, Morse, and URL-encoding are \emph{encodings}: reversible reformatting with no secret key, meant for compatibility or transport, not secrecy. Anyone can reverse them; there is nothing to know. Calling Base64 ``encryption'' is a beginner's error with real consequences -- people have shipped ``encrypted'' configuration files that were merely Base64-encoded, readable by anyone in one command. If there is no key, there is no secret, and there is no cryptography.

\textbf{Cryptography is not a classical ``code.''} Historically a \emph{code} replaced whole words or phrases via a fixed codebook (``BLUEBIRD'' means ``attack at dawn''). Codes operate on meaning; \emph{ciphers} operate mechanically on symbols (letters, bits) regardless of meaning. Modern cryptography is entirely about ciphers and their mathematical generalizations -- the word ``code'' survives mostly in fossils like ``error-correcting code'' and ``decode,'' which are unrelated. When we say ``code'' in this book we mean source code, not a secret codebook.

\begin{trapbox}
``It's encrypted'' is a claim people make about anything they cannot read at a glance -- Base64 blobs, compressed data, binary formats, obfuscated source. None of those involve a key or a security argument, so none of them are encryption. The test is simple: \emph{is there a secret key, and is there a stated adversary the scheme is designed to defeat?} No key and no adversary means no cryptography, however scrambled the output looks.
\end{trapbox}

\section{A Tiny Cipher, and Why It Already Fails}
\label{ch01:tiny-cipher}

Abstractions land better after you have held a concrete example, even a bad one. So before the formal machinery, let us encrypt something by hand, see what ``secure'' would have to mean, and watch exactly how the scheme breaks.

Take the oldest cipher there is, the \emph{shift cipher} (Caesar's cipher). Fix the alphabet $A=0, B=1, \dots, Z=25$ and work modulo $26$. The key is a single number $k \in \{0,1,\dots,25\}$. To encrypt, shift every letter forward by $k$; to decrypt, shift back:
\[
\Enc_k(m) = (m + k) \bmod 26, \qquad \Dec_k(c) = (c - k) \bmod 26,
\]
applied letter by letter. With $k = 3$, the word \texttt{HELLO} becomes:
\[
\texttt{H}(7)\!\to\!10\,\texttt{K},\ \ \texttt{E}(4)\!\to\!7\,\texttt{H},\ \ \texttt{L}(11)\!\to\!14\,\texttt{O},\ \ \texttt{L}\!\to\!\texttt{O},\ \ \texttt{O}(14)\!\to\!17\,\texttt{R},
\]
so the ciphertext is \texttt{KHOOR}. Decryption shifts each letter back by $3$ and recovers \texttt{HELLO}.

\begin{figure}[H]
\centering
\begin{tikzpicture}[font=\small,>=stealth,scale=0.92]
  % Left side: Concentric Cipher Wheel
  \begin{scope}[shift={(0,0)}]
    \draw[thick,sourceblue!50,fill=sourceblue!5] (0,0) circle (2.55cm);
    \draw[thick,targetgreen!50,fill=targetgreen!8] (0,0) circle (1.85cm);
    \draw[thick,structurepurple!40,fill=white] (0,0) circle (1.15cm);
    
    \node[font=\bfseries\scriptsize,structurepurple] at (0,0.3) {Shift $k=3$};
    \draw[->,thick,structurepurple] (-0.4,-0.15) arc (160:380:0.45);

    % Outer letters (Plaintext A-Z) and Inner letters (Ciphertext D-C)
    \foreach \i/\L in {0/A, 1/B, 2/C, 3/D, 4/E, 5/F, 6/G, 7/H, 8/I, 9/J, 10/K, 11/L, 12/M, 13/N, 14/O, 15/P, 16/Q, 17/R, 18/S, 19/T, 20/U, 21/V, 22/W, 23/X, 24/Y, 25/Z} {
      \pgfmathsetmacro{\ang}{90 - \i * (360/26)}
      \node[font=\tiny\bfseries,sourceblue!90!black] at (\ang:2.2) {\L};
    }
    \foreach \i/\L in {0/D, 1/E, 2/F, 3/G, 4/H, 5/I, 6/J, 7/K, 8/L, 9/M, 10/N, 11/O, 12/P, 13/Q, 14/R, 15/S, 16/T, 17/U, 18/V, 19/W, 20/X, 21/Y, 22/Z, 23/A, 24/B, 25/C} {
      \pgfmathsetmacro{\ang}{90 - \i * (360/26)}
      \node[font=\tiny\bfseries,targetgreen!90!black] at (\ang:1.5) {\L};
    }
    \node[above,font=\bfseries\footnotesize,sourceblue] at (0,2.65) {Outer: Plaintext $m$};
    \node[below,font=\bfseries\footnotesize,targetgreen] at (0,-2.65) {Inner: Ciphertext $c = m+3$};
  \end{scope}

  % Right side: Letter-by-letter index transformation
  \begin{scope}[shift={(6.4,0)}]
    \node[above,font=\bfseries\footnotesize] at (1.8,2.45) {Transformation for \texttt{HELLO} $\to$ \texttt{KHOOR}};

    \foreach \k/\x/\char/\idx in {0/0/H/7, 1/0.9/E/4, 2/1.8/L/11, 3/2.7/L/11, 4/3.6/O/14} {
      \node[draw=sourceblue!80,rounded corners=2pt,fill=sourceblue!12,thick,text=sourceblue!90!black,minimum width=0.68cm,minimum height=0.68cm,font=\bfseries] (pt\k) at (\x,1.65) {\char};
      \node[below=1pt,font=\tiny\bfseries,text=gray!75!black] (pti\k) at (pt\k.south) {\idx};
    }
    \node[left,font=\scriptsize\bfseries,sourceblue] at (-0.5,1.65) {Plain $m$:};

    \foreach \k/\x/\cchar/\cidx in {0/0/K/10, 1/0.9/H/7, 2/1.8/O/14, 3/2.7/O/14, 4/3.6/R/17} {
      \node[draw=targetgreen!80,rounded corners=2pt,fill=targetgreen!15,thick,text=targetgreen!90!black,minimum width=0.68cm,minimum height=0.68cm,font=\bfseries] (ct\k) at (\x,-1.25) {\cchar};
      \node[above=1pt,font=\tiny\bfseries,text=gray!75!black] (cti\k) at (ct\k.north) {\cidx};
    }
    \node[left,font=\scriptsize\bfseries,targetgreen] at (-0.5,-1.25) {Cipher $c$:};

    % Arrows connecting index labels cleanly with side label
    \foreach \k in {0, 1, 2, 3, 4} {
      \draw[->,thick,quotientorange] (pti\k.south) -- node[right=1pt,font=\tiny\bfseries,text=quotientorange]{$+3$} (cti\k.north);
    }
    \node[below,font=\scriptsize,align=center] at (1.8,-1.85) {Each index shifts deterministically by $k=3 \pmod{26}$.\\No entropy, no diffusion: word shape is preserved.};
  \end{scope}
\end{tikzpicture}
\caption{The Caesar shift cipher. Left: Concentric wheels shifting the alphabet by $k=3$. Right: Direct letter-to-letter index mapping showing why structural properties pass straight through to the ciphertext.}
\label{ch01:fig-caesar-wheel}
\end{figure}

The scheme is \emph{correct} -- decryption always inverts encryption -- which is the first thing any cipher must get right but says nothing yet about security.

\begin{walkbox}
Before reading on, ask yourself: how many keys are there? If you intercepted \texttt{KHOOR} and knew only that a shift cipher was used (Kerckhoffs: assume the attacker knows the algorithm), how long would it take you to read it?
\end{walkbox}

There are exactly $26$ keys -- and $k=0$ does nothing, so really $25$ useful ones. An attacker simply tries all of them. This is \emph{exhaustive key search}, the most basic attack, and against a $26$-key cipher it is instant.

\begin{attackbox}
\textbf{Breaking the shift cipher two ways.} Intercept \texttt{KHOOR}.
\emph{Brute force:} shift back by $1,2,3,\dots$ and read off the only result that is English. Shifting back by $3$ gives \texttt{HELLO}; the other $24$ candidates are gibberish like \texttt{JGNNQ} or \texttt{GDKKN}. Twenty-five trials, done by hand in under a minute.
\emph{Frequency analysis:} for longer texts you do not even need to try keys. English letter frequencies are lopsided -- \texttt{E} is about $12.7\%$ of letters, \texttt{T} about $9\%$, while \texttt{Z} is near $0.07\%$. A shift cipher slides the whole frequency histogram sideways by $k$ without changing its shape. Find the most common ciphertext letter, assume it is \texttt{E}, and the gap to \texttt{E} reveals $k$. The cipher's fatal flaw is that it \emph{preserves structure}: the statistical fingerprint of the language passes straight through.
\end{attackbox}

The shift cipher is named for Julius Caesar, who according to Suetonius shifted his correspondence by three. It was genuinely useful in an era when most adversaries were illiterate and the very idea of systematic cryptanalysis did not exist. Frequency analysis -- the technique that kills it -- was first described by the polymath al-Kindi in 9th-century Baghdad, in a manuscript on deciphering encrypted messages. That single idea broke essentially every cipher built on simple letter substitution for the next thousand years, and it is why ``it scrambles the letters'' is never, by itself, security (as we will see in Chapter~2).

You might try to patch the shift cipher by allowing an \emph{arbitrary} permutation of the alphabet instead of a mere shift -- the \emph{substitution cipher}. Now the key is a full scrambling of $26$ letters, so there are $26! \approx 4 \times 10^{26}$ keys, a number near $2^{88}$. Surely \emph{that} is unbreakable -- $2^{88}$ is roughly three hundred septillion keys, far too many to brute force by hand or even by machine.

\begin{trapbox}
\textbf{A huge keyspace does not imply security.} The substitution cipher has $\approx 2^{88}$ keys and is broken by a bored teenager in an afternoon -- it is the staple of newspaper cryptograms. Why? A permutation, like a shift, maps each plaintext letter to a \emph{fixed} ciphertext letter, so it still preserves the language's frequency fingerprint: the most common ciphertext symbol is almost surely \texttt{E}, the commonest digraph is probably \texttt{TH}, and so on. Frequency analysis cracks it without ever touching the $2^{88}$ keys. \emph{The size of the keyspace is necessary for security but nowhere near sufficient.} What matters is whether the ciphertext leaks structure -- and these classical ciphers leak it all. We will spend much of the book building ciphers whose output is statistically indistinguishable from random noise, precisely so that no such fingerprint survives.
\end{trapbox}

The shift cipher taught us correctness and exhaustive search; the substitution cipher taught us that keyspace size is not the whole story. Both taught us that ``secure'' is going to need a careful definition, because our intuition (``billions of keys, must be safe!'') just lied to us. Hold these failures in mind; they motivate everything that follows.

\section{The Cast: Alice, Bob, Eve, and Mallory}
\label{ch01:cast}

Cryptographers tell stories about a recurring cast, and the names are not just whimsy -- each one encodes a specific assumption about who is doing what. Getting the cast right is the first step in stating a problem precisely.

\begin{description}
  \item[Alice and Bob] are the two parties who want to communicate securely. By convention Alice initiates and Bob responds, but they are symmetric in spirit. They are the ``good guys,'' though as we will see, even they do not fully trust each other when non-repudiation is at stake.
  \item[Eve] is the \emph{eavesdropper}. She is \emph{passive}: she can see everything that crosses the channel, copy it, store it forever, and analyze it at leisure -- but she does not alter, inject, or delete anything. Eve models wiretaps, packet sniffers, and the intelligence agency recording all traffic for later. Defeating Eve is the confidentiality problem.
  \item[Mallory] is the \emph{active} attacker, the adversary in the middle. Mallory can do everything Eve can, \emph{plus} modify messages, drop them, replay old ones, inject new ones, and impersonate either party. Mallory is far more capable than Eve, and defeating Mallory requires integrity and authentication, not just encryption. The name is a mnemonic: \textbf{Mal}icious, \textbf{Mal}leable, \textbf{M}an-in-the-\textbf{m}iddle.
  \item[Trent] is a \emph{trusted third party} -- an arbiter everyone agrees to believe, such as a certificate authority. Trent is convenient but dangerous: every trusted party is also a single point of failure and a juicy target. Good design minimizes how much you must trust Trent.
\end{description}

You will also meet \textbf{Carol} and \textbf{Dave} (additional legitimate parties when two are not enough), \textbf{Peggy} the \emph{prover} and \textbf{Victor} the \emph{verifier} (in zero-knowledge proofs, where Peggy proves she knows a secret without revealing it), and \textbf{Sybil} (an attacker who forges \emph{many} fake identities at once). The names are a shorthand the whole field shares.

\begin{figure}[H]
\centering
\begin{tikzpicture}[font=\small,>=stealth]
  % Honest communication endpoints
  \node[draw=aliceblue!80,circle,fill=aliceblue!15,minimum size=1.25cm,thick] (A) at (0,0) {\bfseries Alice};
  \node[draw=bobgreen!80,circle,fill=bobgreen!15,minimum size=1.25cm,thick] (B) at (8,0) {\bfseries Bob};

  % Direct honest channel
  \draw[->,thick,sourceblue!80] (A) -- node[above,font=\footnotesize\bfseries,sourceblue]{honest transmission: $c = \Enc_k(m)$} (B);

  % Mallory (Active MitM)
  \node[draw=trapred!80,rounded corners=4pt,fill=trapred!12,minimum width=2.6cm,minimum height=1.0cm,thick] (M) at (4,1.8) {\bfseries Mallory};
  \node[font=\tiny,trapred!90!black,above] at (M.north) {Active Attacker (intercepts, alters, forges)};

  \draw[->,thick,dashed,trapred] (1.3,0.3) to[out=60,in=180] node[left,font=\tiny,trapred]{intercepts $c$} (M.west);
  \draw[->,thick,dashed,trapred] (M.east) to[out=0,in=120] node[right,font=\tiny,trapred]{injects $c'$} (6.7,0.3);

  % Eve (Passive Eavesdropper)
  \node[draw=evered!80,rounded corners=4pt,fill=evered!10,minimum width=2.6cm,minimum height=1.0cm,thick] (E) at (4,-1.8) {\bfseries Eve};
  \node[font=\tiny,evered!90!black,below] at (E.south) {Passive Attacker (observes and records wire)};

  \draw[->,thick,densely dotted,evered] (4,-0.15) -- node[right,font=\tiny,evered]{wiretap: copies $c$ verbatim} (E.north);
\end{tikzpicture}
\caption{The cryptographic threat models. Eve passively records traffic on the channel without disrupting it; confidentiality protects against Eve. Mallory sits directly within the network flow, altering or injecting arbitrary messages; integrity and authentication are required to defeat Mallory.}
\label{ch01:fig-cast}
\end{figure}

\begin{focusbox}
The single most important question to ask about any attacker is \emph{passive or active?} A scheme can be perfectly confidential against Eve and yet immediately broken by Mallory, because confidentiality says nothing about a message you did not even encrypt -- one Mallory injected himself. Most real-world catastrophes (downgrade attacks, man-in-the-middle interception, replay) come from designers who thought about Eve and forgot about Mallory.
\end{focusbox}

\section{Kerckhoffs's Principle: No Security Through Obscurity}
\label{ch01:kerckhoffs}

Here is the foundational rule, and it is counterintuitive enough that beginners resist it: \emph{the security of a cryptographic system must depend only on the secrecy of the key, never on the secrecy of the algorithm.} You should be able to publish your entire design -- the cipher, the protocol, the source code, the exact parameters -- hand it to your enemy, and remain secure as long as the key stays secret.

Auguste Kerckhoffs, a Dutch linguist, laid out six design rules for military ciphers in an 1883 article in the \emph{Journal des sciences militaires}. The second rule is the one that immortalized his name: a system ``must not require secrecy, and it must be able to fall into the hands of the enemy without inconvenience.'' Claude Shannon, founding the mathematical theory of secrecy in 1949, restated it crisply as ``the enemy knows the system'' -- now called \emph{Shannon's maxim}. The principle predates computers by decades, but every modern standard (AES, RSA, the TLS protocol) is fully public precisely because of it.

Why insist on this? Several hard-won reasons.

\textbf{Keys can be changed; algorithms cannot.} If your secret leaks and it was a key, you generate a new key and move on in seconds. If your secret was the algorithm and it leaks, you must redesign, re-implement, re-deploy, and re-audit everything, everywhere, often across hardware you cannot recall. Secrets that are expensive to replace are secrets you will fail to replace.

\textbf{Secret algorithms do not get reviewed.} A cipher's strength is the absence of a clever attack. The only way to gain confidence that no clever attack exists is to let thousands of clever people try to find one. A design kept secret is reviewed by, at most, the handful of people who built it -- the very people least able to see its blind spots. Public algorithms like AES survived years of open cryptanalysis by the world's best before standardization. That public beating is a feature.

\textbf{Obscurity always leaks.} Algorithms get reverse-engineered from binaries, leaked by insiders, extracted from hardware, or simply guessed. Designing as if the algorithm is secret means designing as if a temporary, fragile condition were permanent. It never is.

\begin{attackbox}
\textbf{A5/1 and the GSM phone network.} Early mobile phone encryption (the A5/1 cipher protecting 2G GSM voice calls) was designed in secret in the late 1980s and never published. Secrecy was treated as part of the defense. The design leaked and was reverse-engineered by 1999, and once cryptographers could see it, fatal weaknesses fell out quickly: by the late 2000s, attackers with precomputed tables (``rainbow tables'') could recover the session key for a call in seconds on commodity hardware. Decades of calls had been protected by a cipher whose only real defense -- obscurity -- evaporated the moment someone looked. The contrast with AES, which gained resilience and trust through years of open scrutiny, illustrates the point.
\end{attackbox}

\begin{trapbox}
Kerckhoffs's principle is \emph{not} ``you must publish your source code'' or ``secrecy never helps.'' Keeping deployment details quiet can raise the bar for an attacker as a bonus layer (defense in depth). The principle is narrower and firmer: you may never \emph{rely} on that secrecy for your security. Your scheme must remain secure under the assumption that the attacker knows everything except the key. If revealing your design breaks your system, your system was already broken; you just had not met the right adversary yet.
\end{trapbox}

\paragraph{Enigma: a large keyspace, defeated by structure and operators.} The German military's Enigma machine of World War II had a daily key drawn from a space of around $10^{23}$ possibilities ($\approx 2^{77}$) -- far beyond brute force with the technology of the day. It was broken anyway, first by Polish mathematicians (notably Marian Rejewski) and then on an industrial scale at Britain's Bletchley Park, where Alan Turing's team built electromechanical \emph{bombes} to mechanize the search. The keyspace was never the weak point. The breaks came from structural flaws (no letter ever encrypted to itself -- a constraint the attackers exploited), procedural mistakes (repeated message indicators, predictable plaintext like weather reports and military salutes), and operator laziness. The pattern is the same one the shift cipher showed in miniature and that recurs throughout this book: a giant keyspace is undone by leaked \emph{structure} and human \emph{procedure}, not by counting keys. It is also a clear case for Kerckhoffs -- the Allies effectively obtained the system, and security then rested entirely on the daily key, which the flaws let them recover.

\section{Threat Models: Secure Against Whom, Given What?}
\label{ch01:threat-models}

``Is it secure?'' is not a well-formed question. Secure against a teenager with a laptop, or against a nation-state with a building full of supercomputers and a court order? Secure against someone who only listens, or someone who can also choose which messages you encrypt? A \emph{threat model} is the precise specification of the adversary you are defending against: their goals, their capabilities, and their knowledge. Without it, ``secure'' is marketing, not engineering.

A threat model answers three questions.

\begin{enumerate}
  \item \textbf{What is the attacker trying to do?} (Their \emph{goal}.) Read one specific message? Recover the key? Distinguish two possible plaintexts? Forge a single valid signature? Weaker goals make stronger attacks -- defending against ``learns even one bit of the plaintext'' is harder, hence better, than defending against ``recovers the whole plaintext.''
  \item \textbf{What can the attacker do?} (Their \emph{capabilities}.) This is where we rank attack models by how much access the adversary has to the cryptosystem in action:
    \begin{itemize}
      \item \emph{Ciphertext-only:} the attacker sees ciphertexts and nothing else. The weakest attacker; surviving this is the bare minimum.
      \item \emph{Known-plaintext:} the attacker also sees some plaintext/ciphertext pairs they did not choose (e.g.\ a message with a predictable greeting).
      \item \emph{Chosen-plaintext (CPA):} the attacker can get plaintexts \emph{of their choosing} encrypted. This sounds generous, but it models reality -- an attacker can often trick a system into encrypting data they pick (think of a web server that encrypts whatever you type).
      \item \emph{Chosen-ciphertext (CCA):} the strongest standard model -- the attacker can also get ciphertexts \emph{of their choosing} decrypted (except the one they are trying to break). This models systems that act on decrypted data and leak whether decryption succeeded.
    \end{itemize}
  \item \textbf{What does the attacker know and have?} Computational power, time, side-channel access (timing, power consumption, electromagnetic leakage), and -- per Kerckhoffs -- the full algorithm. We almost always assume the attacker is \emph{computationally bounded}: realistically limited to, say, $2^{100}$ operations, vast but finite. The attacker who can do $2^{256}$ operations can break almost anything by brute force, but such an attacker does not exist and never will (we will do the physics in a moment).
\end{enumerate}

\begin{intuitionbox}
Always design for a \emph{stronger} attacker than you think you face. If your scheme resists chosen-ciphertext attacks, it certainly resists mere eavesdropping; the converse fails badly. The history of broken cryptography is largely the history of designers who modeled a polite attacker and met a rude one. ``We never expected the adversary to be able to do \emph{that}'' is the epitaph on most dead cryptosystems.
\end{intuitionbox}

\begin{connectionbox}
The phrase ``computationally bounded'' is doing heavy lifting and will be made precise using the language of the companion algorithms book. We treat ``feasible'' as ``runs in time polynomial in the key length'' and ``infeasible'' as ``requires super-polynomial (typically exponential) time.'' That is why a cryptographer cares enormously whether the best known attack on a problem is $O(n^3)$ or $2^{n}$: the first is a tractable computation, the second is a wall. Modular exponentiation -- the computational engine of RSA and Diffie-Hellman -- costs $O(\log e)$ multiplications via repeated squaring, which is why public-key crypto is feasible at all. The algorithms companion covers fast exponentiation; the probability companion covers the birthday bound that governs how big a hash must be. We will lean on both.
\end{connectionbox}

\subsection*{How big is $2^{128}$, really?}

Cryptographers throw around numbers like $2^{128}$ as ``infeasible'' and it is worth feeling why. Suppose you build a machine that tries one key every nanosecond -- a billion keys per second. To try all $2^{128}$ keys:

\[
\frac{2^{128}}{10^{9}\ \text{keys/sec}} \approx \frac{3.4 \times 10^{38}}{10^{9}} = 3.4 \times 10^{29}\ \text{seconds}.
\]

There are about $3.15 \times 10^{7}$ seconds in a year, so that is roughly $1.1 \times 10^{22}$ years -- about a trillion times the current age of the universe ($\sim 1.4 \times 10^{10}$ years). Now run a billion such machines in parallel: you save nine orders of magnitude and still need $10^{13}$ years. There is also a hard physical floor: Landauer's principle and the thermodynamics of computation imply that merely \emph{counting} through $2^{128}$ states would consume more energy than is reasonably harvestable. A 128-bit key is not ``hard to brute force.'' It is, against brute force, simply over.

\begin{focusbox}
This is the central bargain of modern cryptography: we make the \emph{legitimate} operation cheap (a fraction of a microsecond to encrypt) and the \emph{attacking} operation astronomically expensive (more than the universe affords), purely by the gap between, say, $2^{30}$ defender work and $2^{128}$ attacker work. Every scheme in this book is an attempt to engineer such a gap and to \emph{prove} the gap really exists.
\end{focusbox}

\section{Randomness: The Resource Everything Rests On}
\label{ch01:randomness}

Hidden inside every secure scheme is an assumption so quiet that beginners miss it entirely: cryptography consumes \emph{genuine unpredictability}. Keys must be unguessable, which means they must be \emph{random} -- chosen by a process the attacker cannot predict or reproduce. The whole edifice rests here. A perfect algorithm fed a predictable ``random'' key is no safer than no algorithm at all, because the attacker simply predicts the key.

Quantify it. If you generate a $128$-bit key by sampling $128$ uniform, independent bits, the attacker faces $2^{128}$ equally likely candidates -- the astronomical wall of \S\ref{ch01:threat-models}. But suppose your random-number generator has an unnoticed flaw and only ever produces keys from a pool of, say, $2^{20}$ possibilities (a million). The key is still $128$ bits \emph{long}; it just is not $128$ bits \emph{unpredictable}. The attacker who knows the flaw -- and per Kerckhoffs we assume they do -- searches only $2^{20} \approx 10^{6}$ keys, a fraction of a second's work. Your $128$-bit cipher now offers $20$ bits of real protection. The length of a key is a promise; the \emph{entropy} behind it is what the promise is worth.

\begin{intuitionbox}
Think of randomness as the fuel and the algorithm as the engine. A flawless engine with contaminated fuel goes nowhere. Most catastrophic real-world breaks are not broken \emph{algorithms} -- AES and RSA stand unbroken -- they are broken \emph{randomness}: predictable keys, reused nonces, weak seeds. The math was perfect; the dice were loaded. When you audit a system, follow the randomness first.
\end{intuitionbox}

\begin{attackbox}
\textbf{When the dice are loaded.} In 2008 a Debian Linux maintainer, cleaning up a memory-handling warning, accidentally removed a line that fed entropy into the OpenSSL random-number generator. For nearly two years, every cryptographic key generated on affected systems was drawn from a pool of only about $2^{15}$ possibilities -- determined essentially by the process ID. Keys that looked like full-strength $2048$-bit RSA or strong SSH keys were, in fact, enumerable by an attacker in minutes. Millions of keys across the internet had to be revoked and regenerated. No cipher was broken. The randomness was. The same theme recurs in embedded devices that generate keys at first boot before they have gathered any entropy, producing predictable or even \emph{shared} keys across thousands of distinct devices -- researchers have factored RSA keys from routers simply because two devices independently generated overlapping ``random'' primes.
\end{attackbox}

This is why cryptographers distinguish a \emph{true random number generator} (TRNG), which harvests physical entropy -- thermal noise, timing jitter, radioactive decay -- from a \emph{cryptographically secure pseudo-random number generator} (CSPRNG), which stretches a short truly-random seed into a long stream that no efficient adversary can distinguish from true randomness. Practical systems gather a little real entropy into a seed, then expand it with a CSPRNG. The catch is the seed: stretch a guessable seed and you get a guessable, infinitely long ``random'' stream. The Debian bug was precisely a poisoned seed.

\begin{trapbox}
The ordinary random functions in most programming languages -- C's \texttt{rand()}, Python's \texttt{random} module, JavaScript's \texttt{Math.random()} -- are \emph{not} cryptographically secure. They are built for simulations and games: fast, statistically smooth, and \emph{predictable}, because an attacker who sees a few outputs can recover the internal state and predict all future ones. Never generate keys, nonces, or tokens with them. Use the operating system's cryptographic source (\texttt{/dev/urandom}, \texttt{getrandom()}, Python's \texttt{secrets} or \texttt{os.urandom}, Java's \texttt{SecureRandom}). The \texttt{os.urandom} call in the listing of \S\ref{ch01:what-is-secure} was deliberate -- it is the right source -- even though that example's \emph{usage} of the output was wrong.
\end{trapbox}

\begin{connectionbox}
Whether a stretched seed ``looks random'' to every efficient adversary is the formal notion of a \emph{pseudo-random generator}, and proving a CSPRNG secure reduces -- as always -- to some hardness assumption. The probability companion provides the statistical-distance and indistinguishability tools used to state ``cannot be told apart from uniform''; the number-theory chapters provide candidate hard problems (like the difficulty of certain residue questions) from which provably secure generators can be built. Randomness is not a footnote to cryptography; it is half of it.
\end{connectionbox}

\section{Symmetric vs.\ Asymmetric: A First Look}
\label{ch01:sym-vs-asym}

There are two great families of cryptographic schemes, distinguished by how keys are arranged. We preview them now and devote later chapters to each.

Before naming the families, let us write down what an encryption scheme even \emph{is}, as a piece of syntax. This three-algorithm template recurs for every scheme in the book, so it is worth fixing now.

\begin{definition}[Symmetric encryption scheme]
\label{ch01:def-scheme}
A symmetric (private-key) encryption scheme is a triple of algorithms $(\Gen, \Enc, \Dec)$ over a message space $\mathcal{M}$, key space $\mathcal{K}$, and ciphertext space $\mathcal{C}$:
\begin{itemize}
  \item $\Gen$ outputs a key $k \in \mathcal{K}$, typically by sampling uniformly at random: $k \getsr \mathcal{K}$.
  \item $\Enc$ takes a key $k$ and message $m \in \mathcal{M}$ and outputs a ciphertext $c \gets \Enc_k(m)$. It may be randomized -- and for serious security it \emph{must} be.
  \item $\Dec$ takes a key $k$ and ciphertext $c$ and outputs a message $m \gets \Dec_k(c)$.
\end{itemize}
The scheme is \emph{correct} if decryption inverts encryption: for every key $k$ in the support of $\Gen$ and every $m \in \mathcal{M}$,
\[
\Dec_k(\Enc_k(m)) = m.
\]
\end{definition}

Correctness is the part beginners never forget; security is the part the definition above says \emph{nothing} about. The shift cipher of \S\ref{ch01:tiny-cipher} fits this template perfectly and is correct -- and useless. Syntax is not security. With the template in hand, the two families differ in just one place: how the key for $\Enc$ relates to the key for $\Dec$.

\textbf{Symmetric (secret-key) cryptography.} Alice and Bob share one secret key $k$. The same $k$ encrypts and decrypts:
\[
c = \Enc_k(m), \qquad m = \Dec_k(c).
\]
Symmetric schemes are fast -- modern CPUs encrypt gigabytes per second -- and they are how the bulk of real data gets protected. AES is the dominant example. The catch is the \emph{key distribution problem}: before they can talk, Alice and Bob must already share $k$ through some secure channel. If they could do that reliably, why not just send the message that way? For two people who can meet in person once, fine. For a browser connecting to a website it has never contacted, hopeless.

\textbf{Asymmetric (public-key) cryptography.} Each party has a \emph{key pair}: a public key $pk$ that anyone may know, and a private key $sk$ kept secret. The keys are mathematically linked so that what one does, only the other can undo. For encryption: anyone encrypts to Bob using his public key, and only Bob, holding $sk$, can decrypt:
\[
c = \Enc_{pk_B}(m), \qquad m = \Dec_{sk_B}(c).
\]
For signatures the roles flip: Alice signs with her \emph{private} key, and anyone can verify with her public key -- giving authentication and non-repudiation in one stroke. Public-key crypto (RSA, Diffie-Hellman, elliptic curves) solves key distribution: Bob publishes $pk_B$ to the world, and strangers can send him secrets with no prior shared secret. The cost is speed -- asymmetric operations are orders of magnitude slower than symmetric ones.

\begin{intuitionbox}
\emph{Symmetric} is one key shared by two trusting parties -- like a single physical key both copies fit the same lock. \emph{Asymmetric} is a mailbox: anyone can drop a letter through the public slot (the public key), but only the holder of the physical key (the private key) can open it and take the letters out. The mailbox slot can be photographed and broadcast on television; it does not help the thief open the box.
\end{intuitionbox}

In practice, real systems are \emph{hybrid}. Because public-key operations are slow, you use them only to do the hard part -- agree on a fresh symmetric key over an insecure channel -- and then switch to fast symmetric encryption for the actual data. Every HTTPS connection works this way: a public-key handshake establishes a shared secret, then AES carries the gigabytes. We get the key-distribution win of asymmetric crypto and the speed of symmetric crypto in one design.

\begin{figure}[H]
\centering
\begin{tikzpicture}[font=\small,>=stealth]
  % Symmetric Panel
  \begin{scope}[shift={(0,0)}]
    \draw[thick,sourceblue!40,fill=sourceblue!5,rounded corners=6pt] (-0.3,-1.3) rectangle (5.5,2.7);
    \node[font=\bfseries\normalsize,sourceblue!90!black] at (2.6,2.3) {Symmetric Cryptography};

    \node[draw=aliceblue!80,circle,fill=aliceblue!20,thick,minimum size=0.95cm] (sa) at (0.6,0.75) {\scriptsize Alice};
    \node[draw=bobgreen!80,circle,fill=bobgreen!20,thick,minimum size=0.95cm] (sb) at (4.6,0.75) {\scriptsize Bob};

    % Key token
    \node[draw=quotientorange!80,fill=quotientorange!15,rounded corners=2pt,font=\tiny\bfseries,inner sep=2.5pt] (skey) at (2.6,1.65) {Shared Key $k$};
    \draw[dashed,quotientorange!80] (sa.north east) -- (skey.west);
    \draw[dashed,quotientorange!80] (skey.east) -- (sb.north west);

    \draw[->,thick,sourceblue!80] (sa) -- node[above=2pt,font=\tiny\bfseries]{$c = \Enc_k(m)$} (sb);
    \node[font=\tiny,gray!85!black] at (2.6,0.35) {$m = \Dec_k(c)$};

    \node[font=\tiny,align=center,quotientorange!90!black] at (2.6,-0.65) {
      \textbf{Core bottleneck:} Key distribution.\\
      How do Alice and Bob safely share $k$ beforehand?
    };
  \end{scope}

  % Asymmetric Panel
  \begin{scope}[shift={(6.4,0)}]
    \draw[thick,targetgreen!40,fill=targetgreen!5,rounded corners=6pt] (-0.3,-1.3) rectangle (5.5,2.7);
    \node[font=\bfseries\normalsize,targetgreen!90!black] at (2.6,2.3) {Asymmetric Cryptography};

    \node[draw=aliceblue!80,circle,fill=aliceblue!20,thick,minimum size=0.95cm] (aa) at (0.6,0.75) {\scriptsize Alice};
    \node[draw=bobgreen!80,circle,fill=bobgreen!20,thick,minimum size=0.95cm] (ab) at (4.6,0.75) {\scriptsize Bob};

    % Public/Private key tokens
    \node[draw=sourceblue!80,fill=sourceblue!15,rounded corners=2pt,font=\tiny\bfseries,inner sep=2.5pt] (pubkey) at (1.5,1.65) {Public $pk_B$};
    \node[draw=structurepurple!80,fill=structurepurple!15,rounded corners=2pt,font=\tiny\bfseries,inner sep=2.5pt] (privkey) at (4.3,-0.15) {Private $sk_B$};

    \draw[->,thick,targetgreen!80] (aa) -- node[above=2pt,font=\tiny\bfseries]{$c = \Enc_{pk_B}(m)$} (ab);
    \node[font=\tiny,gray!85!black] at (2.6,0.35) {$m = \Dec_{sk_B}(c)$};

    \node[font=\tiny,align=center,targetgreen!90!black] at (2.6,-0.75) {
      \textbf{Breakthrough:} Zero prior shared secret.\\
      Anyone encrypts with $pk_B$; only $sk_B$ decrypts.
    };
  \end{scope}
\end{tikzpicture}
\caption{Symmetric vs.\ asymmetric architecture. In symmetric schemes, one shared key does both encryption and decryption, creating a chicken-and-egg distribution problem. In asymmetric schemes, the key pair splits the roles: public encryption is open to all, while private decryption remains bound to a single secret.}
\label{ch01:fig-sym-asym}
\end{figure}

\begin{historybox}
For all of recorded history until the 1970s, \emph{all} cryptography was symmetric, and the key-distribution problem was solved by couriers, codebooks, and trust -- expensively and unreliably. The conceptual leap that public keys were even \emph{possible} came in 1976 from Whitfield Diffie and Martin Hellman (building on a suggestion by Ralph Merkle), and shortly after in the RSA scheme of Rivest, Shamir, and Adleman. We later learned that researchers at the British intelligence agency GCHQ (James Ellis, Clifford Cocks, Malcolm Williamson) had discovered essentially the same ideas between 1970 and 1973, but their work stayed classified for over two decades. The public discovery sparked a fight with the US National Security Agency, which tried to restrict open academic publication under export controls; when open research prevailed, it made web commerce, TLS, and secure messaging possible, and it is itself a small monument to Kerckhoffs. Today the main open challenge for public-key cryptography is post-quantum migration: because Shor's algorithm on a quantum computer can factor integers and compute discrete logarithms in polynomial time, standard systems are moving toward lattice-based hardness assumptions.
\end{historybox}

\section{What Does ``Secure'' Even Mean?}
\label{ch01:what-is-secure}

Now the deep part. We have used the word ``secure'' a dozen times; it is time to admit it is treacherous. The history of cryptography is a graveyard of schemes that everyone, including experts, \emph{felt} were secure -- until they weren't. Intuition is not enough. The entire modern field exists to replace ``I can't see how to break it'' with ``I can prove no efficient attacker can break it, or breaking it would solve a problem the whole world believes is hard.''

\subsection*{Why intuition fails}

Consider a tempting argument. ``I XOR my message with a long pseudo-random keystream. The keystream looks completely random, so the ciphertext looks completely random, so the attacker learns nothing.'' That feels airtight. But it quietly assumes the attacker only ever sees \emph{one} ciphertext and never sees any plaintext. Reuse the keystream on two messages $m_1, m_2$, and the attacker computes
\[
c_1 \xor c_2 = (m_1 \xor k) \xor (m_2 \xor k) = m_1 \xor m_2,
\]
and the key has vanished, leaving the two plaintexts XORed together -- often enough to recover both. The intuition was not wrong, exactly; it answered the wrong threat model. This exact mistake -- keystream reuse -- has broken real systems repeatedly (the WEP WiFi protocol, the wartime Soviet one-time-pad reuse that the VENONA project exploited, Microsoft's PPTP). Feeling secure is not being secure.

\begin{trapbox}
``It looks random'' is not a security argument. ``I tried for an afternoon and couldn't break it'' is not a security argument. ``It's been used for years'' is not a security argument (A5/1 was). The only thing that counts is a precise claim about a precise adversary, backed by either a proof or, at minimum, sustained public failure to break it despite strong incentives. Anything less is hope.
\end{trapbox}

\subsection*{The three ingredients of a security definition}

A modern security definition always specifies three things together. Drop any one and the claim becomes meaningless.

\begin{enumerate}
  \item \textbf{A security goal} -- exactly what the attacker must fail to do. For encryption the gold standard is \emph{indistinguishability}: the attacker, having chosen two equal-length messages $m_0, m_1$, cannot tell which one a given ciphertext encrypts, any better than flipping a coin. This is far stronger than ``cannot recover the plaintext,'' and deliberately so: it forbids leaking even a single bit, even the answer to ``is the salary above \$100k?''
  \item \textbf{A threat model} -- exactly what the attacker may do, from \S\ref{ch01:threat-models}: ciphertext-only, chosen-plaintext, chosen-ciphertext, and how much computation they get.
  \item \textbf{A quantified bound} -- exactly how badly they fail. We do not demand the attacker \emph{never} succeeds; with enough luck anyone can guess a key. We demand their advantage over guessing is \emph{negligible}: it shrinks faster than any inverse polynomial in the key length. Double the key length and the attacker's success chance should crater.
\end{enumerate}

Putting these together: a scheme is \emph{secure} if every efficient (polynomial-time) attacker, operating within the stated threat model, achieves only negligible advantage at the stated goal. That is the shape of every real theorem in this book.

\begin{formalbox}
We will write security as a game. A challenger and an adversary $\Adv$ interact by fixed rules; $\Adv$ ``wins'' if it achieves the goal. The scheme is secure if for every probabilistic polynomial-time ($\PPT$) adversary,
\[
\Pr[\,\Adv \text{ wins}\,] \le \tfrac{1}{2} + \negl(n),
\]
where $n$ is the security parameter (think key length) and $\negl(n)$ is a function smaller than $1/p(n)$ for every polynomial $p$ and all large $n$. The $\tfrac12$ is the trivial coin-flip baseline for a yes/no challenge; the whole content is that no clever strategy beats coin-flipping by more than a vanishing margin. We unfold this machinery carefully in the chapters on perfect and computational secrecy. For now, absorb the \emph{shape}: quantify over all efficient adversaries, bound their advantage by a negligible function.
\end{formalbox}

The indistinguishability game deserves a picture, because its structure is the skeleton of nearly every security proof you will see. The adversary picks two messages it considers maximally telling; the challenger secretly flips a coin $b$, encrypts $m_b$, and hands back the ciphertext; the adversary must guess $b$. If even the cleverest adversary cannot beat a coin flip by a non-negligible margin, the ciphertext is leaking nothing useful about which message it hides.

\begin{figure}[H]
\centering
\begin{tikzpicture}[font=\small,>=stealth]
  % Lanes
  \draw[thick,trapred!30,fill=trapred!5,rounded corners=6pt] (-0.5,-3.9) rectangle (3.5,0.75);
  \draw[thick,bobgreen!30,fill=bobgreen!5,rounded corners=6pt] (5.5,-3.9) rectangle (9.5,0.75);

  % Column Headers
  \node[draw=trapred!80,fill=trapred!20,rounded corners=3pt,font=\bfseries,minimum width=3.4cm,minimum height=0.6cm] at (1.5,0.3) {Adversary $\Adv$};
  \node[draw=bobgreen!80,fill=bobgreen!20,rounded corners=3pt,font=\bfseries,minimum width=3.4cm,minimum height=0.6cm] at (7.5,0.3) {Challenger $\mathcal{C}$};

  % Step 1: Send m0, m1
  \draw[->,thick,sourceblue!90!black] (2.8,-0.45) -- node[above=2pt,font=\scriptsize\bfseries]{(1) $m_0, m_1$ \text{ with } $|m_0|=|m_1|$} (6.2,-0.45);

  % Step 2: Challenger actions
  \node[draw=bobgreen!70,fill=white,rounded corners=2pt,align=center,font=\scriptsize,inner sep=4pt] at (7.5,-1.45) {
    (2) Sample $k \gets \mathcal{K}$\\[3pt]
    Toss coin $b \gets \{0,1\}$\\[3pt]
    Compute $c^* = \Enc_k(m_b)$
  };

  % Step 3: Send c*
  \draw[<-,thick,targetgreen!90!black] (2.8,-2.5) -- node[above=2pt,font=\scriptsize\bfseries]{(3) Challenge ciphertext $c^*$} (6.2,-2.5);

  % Step 4: Adversary guess
  \node[draw=trapred!70,fill=white,rounded corners=2pt,align=center,font=\scriptsize,inner sep=4pt] at (1.5,-3.15) {
    (4) Efficient computation\\[2pt]
    Output guess $b' \in \{0,1\}$
  };

  % Bottom banner: Advantage & Win condition
  \node[draw=structurepurple!50,fill=structurepurple!8,rounded corners=4pt,thick,align=center,font=\footnotesize,minimum width=10.2cm,inner sep=5pt] at (4.5,-4.65) {
    \textbf{Win Condition:} $\Adv$ wins if $b' = b$. \quad
    \textbf{Advantage:} $\Adv_{\Pi}(\Adv) = \left|\Pr[b'=b] - \tfrac{1}{2}\right| \le \negl(n)$
  };
\end{tikzpicture}
\caption{The indistinguishability game (IND-EAV / IND-CPA). The adversary picks two candidate messages of equal length; the challenger encrypts one chosen by a fair coin flip. Security guarantees that no efficient algorithm can guess the coin any better than pure chance.}
\label{ch01:fig-ind-game}
\end{figure}

It pays to be careful about \emph{advantage} versus raw probability, because the difference is where the content lives. In a yes/no game an adversary that ignores the ciphertext and guesses $b' \getsr \{0,1\}$ already wins with probability $\tfrac12$. So $\Pr[b'=b] = \tfrac12$ means ``learned nothing,'' not ``always wrong.'' We therefore measure the \emph{advantage}
\[
\Adv = \abs[\big]{\Pr[b'=b] - \tfrac12},
\]
the margin by which the adversary beats blind guessing. Security demands $\Adv$ be negligible. Concretely, an adversary winning with probability $0.5000001$ has advantage $10^{-7}$ -- tiny, plausibly tolerable. An adversary winning with probability $0.75$ has advantage $0.25$ -- a quarter of the way to certainty, completely broken. And an adversary winning with probability $0.5$ exactly has advantage $0$ -- the ideal. Reading a security theorem is largely reading a bound on this one quantity.

\begin{example}[Why deterministic encryption loses the game]
\label{ch01:ex-deterministic}
Suppose $\Enc$ is \emph{deterministic}: encrypting the same message always yields the same ciphertext (the shift cipher is like this). Then the game is trivial to win. The adversary submits any two distinct messages $m_0 \ne m_1$, and -- if it is allowed even a single chosen-plaintext query, or simply knows the encryptions of $m_0$ and $m_1$ from earlier traffic -- it precomputes $\Enc_k(m_0)$, compares it to the challenge ciphertext $c$, and answers $b'=0$ if they match and $b'=1$ otherwise. It wins with probability $1$, not $\tfrac12$. The advantage is $\tfrac12$, the largest possible -- catastrophically non-negligible. This is the formal reason \emph{secure encryption must be randomized}: the same plaintext encrypted twice must (with overwhelming probability) give different ciphertexts, or an attacker learns when you repeat yourself. ``This message equals that earlier one'' is exactly the kind of bit indistinguishability forbids.
\end{example}

\begin{connectionbox}
The word ``negligible'' has a precise probabilistic meaning developed in the companion probability book: a function $\negl(n)$ that is eventually smaller than every inverse polynomial $1/n^c$. The reason this is the right cutoff is an amplification argument -- if an attacker had even a $1/n^c$ chance per try, repeating a polynomial number of times would boost success to near certainty, so anything non-negligible is, for a patient attacker, eventually fatal. The probability companion also develops the \emph{birthday bound}, which dictates that a hash function needs a $2n$-bit output to resist collision-finding at the $2^n$ level, and the union and Markov bounds we use constantly to push success probabilities around in reductions. Keep that book within reach; every security definition in this book is a quantitative probability statement.
\end{connectionbox}

\subsection*{Two routes to confidence: perfect and computational security}

There are two ways to back a security claim.

\emph{Information-theoretic (perfect) security} promises that the ciphertext reveals \emph{nothing} about the plaintext, even to an adversary with unlimited computing power. Shannon proved the one-time pad achieves this -- but at a brutal price: the key must be as long as the message and never reused. Perfect security is real, provable, and mostly impractical.

\emph{Computational security} -- the foundation of modern practical protocols -- promises something more modest and more useful: the ciphertext might in principle reveal the plaintext, but extracting it would require more computation than any feasible adversary commands. In practical cryptography, these proofs are almost always \emph{conditional}: ``if factoring large integers is hard, then RSA is secure.'' We do not know how to prove factoring is hard outright (that would resolve deep open questions in complexity theory), but we reduce the security of a whole scheme to one clean, ancient, heavily-studied problem. If the scheme breaks, one of a short list of famous problems must have fallen -- and we would all have much bigger news to discuss.

\begin{whybox}
\textbf{Why reductions are the heart of the field.} Suppose I prove: ``any efficient algorithm that breaks my encryption scheme can be transformed into an efficient algorithm that factors large numbers.'' Then breaking my scheme is \emph{at least as hard} as factoring. I have not proven factoring is hard -- nobody has -- but I have concentrated all the risk into a single, well-studied assumption that the world's best mathematicians have failed to crack for centuries. That is enormously more trustworthy than ``I, personally, couldn't break it.'' A reduction converts ``trust my cleverness'' into ``trust the failure of every number theorist since Fermat.'' This is the discipline that separates cryptography from codemaking. Almost every theorem you meet here has the form: \emph{if problem $P$ is hard, then scheme $S$ is secure}, proved by turning an attacker on $S$ into a solver for $P$ (Figure~\ref{ch01:fig-reduction}).
\end{whybox}

\begin{figure}[H]
\centering
\begin{tikzpicture}[font=\small,>=stealth]
  % Outer solver box B
  \draw[thick,sourceblue!70!black,fill=sourceblue!6,rounded corners=6pt] (1.3,-1.65) rectangle (8.7,1.65);
  \node[font=\bfseries\scriptsize,sourceblue!90!black] at (5.0,1.25) {Reduction wrapper $\mathcal{B}$ (solves hard problem $P$)};

  % Inner adversary box A
  \node[draw=trapred!80!black,fill=trapred!12,thick,rounded corners=4pt,align=center,minimum width=3.4cm,minimum height=1.1cm] (adv) at (5.0,-0.35) {
    \textbf{\color{trapred!90!black}Hypothetical $\Adv$}\\
    \scriptsize breaks scheme $S$
  };

  % External input x
  \node[draw=quotientorange!80!black,fill=quotientorange!12,thick,rounded corners=3pt,align=center,inner sep=4pt] (inp) at (-0.9,0) {
    Instance $x$\\
    \scriptsize (e.g.\ $n = pq$)
  };

  % External output y
  \node[draw=targetgreen!80!black,fill=targetgreen!12,thick,rounded corners=3pt,align=center,inner sep=4pt] (out) at (10.9,0) {
    Solution $y$\\
    \scriptsize (e.g.\ factors $p, q$)
  };

  % Arrows
  \draw[->,thick,quotientorange!85!black] (inp.east) -- (1.3,0);
  \draw[->,thick,sourceblue!85!black] (1.9,0.3) -- node[above=2pt,font=\scriptsize]{embeds $x$} (3.3,0.05);
  \draw[->,thick,trapred!85!black] (6.7,0.05) -- node[above=2pt,font=\scriptsize]{outputs break} (8.1,0.3);
  \draw[->,thick,targetgreen!85!black] (8.7,0) -- (out.west);
\end{tikzpicture}
\caption{Proof by reduction. A wrapper algorithm $\mathcal{B}$ takes a hard-problem instance $x$, embeds it into a simulated challenge for a hypothetical attacker $\Adv$ against scheme $S$, and translates $\Adv$'s break into a solution $y$ for $P$. Because $P$ is believed hard, no efficient $\Adv$ can exist.}
\label{ch01:fig-reduction}
\end{figure}

\begin{implementbox}
Here is the keystream-reuse pitfall in a few lines of Python, so you can \emph{see} the key dissolve. This is a demonstration of a \textbf{broken} use of encryption -- the very mistake \S\ref{ch01:what-is-secure} warns against. \textbf{Learning only: never reuse a keystream, and never ship this.}
\begin{lstlisting}[language=Python]
import os

def xor(a, b):
    return bytes(x ^ y for x, y in zip(a, b))

m1 = b"attack at dawn!!"
m2 = b"retreat at noon."   # same length as m1

# THE BUG: one keystream reused for two messages.
keystream = os.urandom(len(m1))
c1 = xor(m1, keystream)
c2 = xor(m2, keystream)

# The attacker never sees the key, yet:
leak = xor(c1, c2)
# leak == m1 XOR m2, completely independent of the key.
assert leak == xor(m1, m2)

# If the attacker guesses (or knows) m1, m2 falls out for free:
recovered_m2 = xor(leak, m1)
assert recovered_m2 == m2
print("recovered:", recovered_m2)   # b"retreat at noon."
\end{lstlisting}
The key was a perfect, unguessable random string. It did not matter. The \emph{usage} violated the threat model the one-time pad requires (each key used once), and the math punished it exactly. Security lives in definitions and usage, not in vibes about randomness.
\end{implementbox}

\section{The Security Mindset}
\label{ch01:mindset}

Everything above converges on a habit of thought. A builder asks ``how do I make this work?'' A cryptographer asks ``how does this fail, and who makes it fail?'' The security mindset is the reflex to look at any system and immediately hunt for the gap between what the designer assumed and what an adversary can actually do.

Concretely, when you meet a scheme, interrogate it:

\begin{itemize}
  \item \textbf{What exactly is the secret, and what happens when it leaks?} If the answer is ``the algorithm,'' stop -- Kerckhoffs already lost this for you.
  \item \textbf{Who is the attacker, precisely?} Passive Eve or active Mallory? What attack model -- ciphertext-only or chosen-ciphertext? Name the strongest realistic adversary, then check the scheme against \emph{that} one, not the polite one.
  \item \textbf{What is the security goal, stated as something the attacker must fail to do?} If you cannot state a crisp win condition for the adversary, you cannot reason about security at all.
  \item \textbf{Where is the proof, or the assumption?} Is the claim ``provably as hard as factoring,'' or merely ``nobody's broken it yet''? Both can be acceptable -- AES has no reduction, only a long unbroken public record -- but you must know which one you are leaning on.
  \item \textbf{What did the model ignore?} The math can be perfect while the implementation leaks the key through timing, power draw, or a careless error message. The most elegant proof says nothing about a side channel it never modeled.
\end{itemize}

\begin{attackbox}
\textbf{Padding-oracle attacks: when the model and the world diverge.} A cipher can be provably secure against a chosen-plaintext attacker and still be torn apart in deployment. The padding-oracle attack targets systems that, on receiving a malformed ciphertext, reveal -- through an error message, a timing difference, or a behavioral quirk -- \emph{whether the decryption padding was valid}. That single leaked bit per query, repeated, lets an attacker decrypt traffic without ever knowing the key. It works because real systems often act as a \emph{chosen-ciphertext oracle} that the original ``CPA-secure'' analysis never accounted for. This is the security mindset in action: the proof assumed one model, reality presented another, and the discrepancy was precisely where the exploit lived. The fix -- authenticated encryption that rejects tampered ciphertexts before they reach the padding check -- closes the gap by enlarging the model to match reality.
\end{attackbox}

\begin{intuitionbox}
A cryptographer is professionally paranoid in a specific, disciplined way. Not ``everything is dangerous'' (that is paralysis) but ``state your assumptions, then assume the adversary violates every one you did not nail down.'' The discipline is what turns paranoia into proofs. You assume the worst about the adversary so that you can make precise, defensible promises about the system -- and then you test those promises by handing your design to the smartest enemies you can find and inviting them to break it.
\end{intuitionbox}

\begin{applicationbox}
The Signal messaging protocol is a fine specimen of this mindset shipped at scale. It does not merely encrypt messages; it is engineered against a remarkably strong attacker. \emph{Forward secrecy}: if your phone is seized and keys extracted today, yesterday's messages stay safe, because each message used a key already deleted. \emph{Post-compromise security} (``future secrecy''): even after a temporary key compromise, the protocol heals and future messages regain protection as fresh key material ratchets in. These are not goals a casual designer would think to chase -- they answer the question ``what if the attacker eventually gets the key anyway?'' That question, asked early and answered in the design, is the security mindset made concrete. We build toward Signal's ratchet in Chapter~27.
\end{applicationbox}

\section{Composing the Goals: A Single Message, Several Guarantees}
\label{ch01:composition}

Real messages usually need more than one guarantee at once, and the order in which you bolt mechanisms together is itself a security decision with a wrong answer. Consider Alice sending Bob a message that must be \emph{both} confidential (Eve learns nothing) \emph{and} integrity-protected (Mallory cannot alter it undetected). She has two primitives: an encryption scheme for confidentiality and a message-authentication code (MAC), a keyed function that produces a short tag $t = \MAC_{k}(x)$ which Bob can recompute to verify $x$ was not tampered with. How should she combine them?

There are three natural orders, and they are not equally safe.

\begin{itemize}
  \item \textbf{Encrypt-and-MAC}: send $\Enc_{k_1}(m)$ alongside $\MAC_{k_2}(m)$. The tag is computed on the \emph{plaintext} and sent in the clear-ish next to the ciphertext. Problem: a MAC guarantees integrity, not secrecy, so the tag may leak information about $m$ -- and since the MAC of a fixed $m$ is fixed, it can betray when Alice repeats a message. Fragile.
  \item \textbf{MAC-then-encrypt}: compute $t = \MAC_{k_2}(m)$, then send $\Enc_{k_1}(m \concat t)$. Better, but the tag protects the plaintext, so Bob must \emph{decrypt first} to check integrity -- meaning he runs his decryption machinery on attacker-controlled ciphertext before he knows it is authentic. That is exactly the opening the padding-oracle attack of \S\ref{ch01:mindset} exploited.
  \item \textbf{Encrypt-then-MAC}: compute $c = \Enc_{k_1}(m)$, then $t = \MAC_{k_2}(c)$, and send $c \concat t$. Bob checks the tag on the \emph{ciphertext} first; if it fails, he discards the message \emph{without ever decrypting}. The attacker cannot get the decryption machinery to run on anything it forged. This order is the one modern practice settled on, and with the right primitives it is provably secure.
\end{itemize}

\begin{figure}[H]
\centering
\begin{tikzpicture}[font=\small,>=stealth]
  % Sender Lane (Alice)
  \begin{scope}[shift={(0,2.0)}]
    \node[font=\bfseries\footnotesize,sourceblue,anchor=east] at (-0.3,0) {Alice (Sender):};
    \node (m) at (0.3,0) {$m$};
    \node[draw=sourceblue!80,fill=sourceblue!12,thick,rounded corners=3pt] (enc) at (1.9,0) {$\Enc_{k_1}$};
    \node (c1) at (3.6,0) {$c$};
    
    % branch for MAC
    \node[draw=targetgreen!80,fill=targetgreen!12,thick,rounded corners=3pt] (mac) at (5.2,-0.8) {$\MAC_{k_2}$};
    \node (t) at (7.0,-0.8) {$t$};

    \draw[->,thick] (m) -- (enc);
    \draw[->,thick] (enc) -- (c1);
    \draw[->,thick] (c1.south) |- (mac.west);
    \draw[->,thick] (mac) -- (t);

    \node[draw=structurepurple!60,fill=structurepurple!10,rounded corners=3pt,thick,inner sep=4pt] (pkt) at (8.5,-0.4) {Wire: $(c, t)$};
    \draw[->,thick] (c1.east) -| (pkt.north);
    \draw[->,thick] (t.east) -| (pkt.south);
  \end{scope}

  % Divider
  \draw[dashed,gray!50] (-3.1,0.7) -- (9.8,0.7);

  % Receiver Lane (Bob)
  \begin{scope}[shift={(0,-0.6)}]
    \node[font=\bfseries\footnotesize,bobgreen,anchor=east] at (-0.3,0) {Bob (Receiver):};
    \node[draw=structurepurple!60,fill=structurepurple!10,rounded corners=3pt,thick,inner sep=4pt] (inpkt) at (0.7,0) {$(c, t)$};

    % Verification Gatekeeper
    \node[draw=targetgreen!80,fill=targetgreen!15,thick,rounded corners=3pt,align=center,font=\scriptsize] (vfy) at (3.3,0) {$\Vrfy_{k_2}(c,t)$};
    \draw[->,thick] (inpkt) -- (vfy);

    % Decision branches
    \node[draw=trapred!80,fill=trapred!15,rounded corners=2pt,font=\tiny\bfseries,align=center] (reject) at (5.7,-0.9) {Reject $\bot$\\ (Drop packet)};
    \draw[->,thick,trapred] (vfy.south) |- node[near end,below=1pt,font=\tiny]{Invalid Tag} (reject.west);

    \node[draw=sourceblue!80,fill=sourceblue!12,thick,rounded corners=3pt] (dec) at (6.0,0.7) {$\Dec_{k_1}$};
    \draw[->,thick,targetgreen!80!black] (vfy.north) |- node[near end,above=1pt,font=\tiny]{Valid Tag} (dec.west);

    \node (rec_m) at (7.8,0.7) {$m$};
    \draw[->,thick] (dec) -- (rec_m);

    \node[align=center,font=\tiny,quotientorange!90!black] at (4.5,-1.85) {
      \textbf{Firewall guarantee:} Tampered ciphertexts are dropped before decryption.\\
      Padding-oracle attacks cannot trigger because invalid payloads never reach $\Dec$.
    };
  \end{scope}
\end{tikzpicture}
\caption{Encrypt-then-MAC (EtM). Tagging the ciphertext creates a cryptographic firewall: Bob verifies authenticity \emph{before} handing the ciphertext to the decryption routine, cutting off padding-oracle attacks at the door.}
\label{ch01:fig-etm}
\end{figure}

The takeaway is not the specific ordering -- we will prove the relevant theorems later, in Chapters~15 and~16. The takeaway is that \emph{combining secure pieces does not automatically yield a secure whole}. Composition has rules, and ignoring them reintroduces exactly the attacks each piece was meant to stop. Modern designs sidestep the footgun by providing a single primitive, \emph{authenticated encryption with associated data} (AEAD) -- algorithms like AES-GCM and ChaCha20-Poly1305 -- that delivers confidentiality and integrity together, in the proven-correct order, behind one interface. When you reach for crypto in real code, reach for an AEAD, not a bare cipher.

\begin{applicationbox}
Every modern TLS connection (TLS 1.3, the version behind today's HTTPS) uses exactly one category of symmetric primitive for record protection: an AEAD cipher, almost always AES-GCM or ChaCha20-Poly1305. The protocol designers removed the older, error-prone MAC-then-encrypt constructions \emph{from the standard entirely}, because a decade of padding-oracle and related attacks showed that leaving the composition to implementers was a reliable source of vulnerabilities. That is the lesson of this section baked directly into the most widely deployed security protocol on Earth: give people the safe composition as the only option, and the whole class of ordering bugs disappears.
\end{applicationbox}

\section{What's Ahead}
\label{ch01:roadmap}

You now have the frame. In Chapter~2 we step through the history of secret writing -- from Caesar and Vigen\`ere through Enigma -- to see why intuitive ciphers collapsed under frequency analysis and structural leaks. After that, Chapters~3 through~6 build the mathematics secrecy is made of (modular arithmetic, primes and Euclid's algorithm, groups, and finite fields), followed by the one-time pad and perfect secrecy, stream and block ciphers (AES), MACs and hash functions, public-key schemes (Diffie--Hellman, RSA, elliptic curves) and signatures, and finally the protocols that run the internet: TLS, WiFi, Signal, and SSH. Keep the four goals, the cast, Kerckhoffs's principle, and the three-part security definition in your pocket. Every chapter ahead is a variation on the questions of this one: secure against whom, given what, proven how?

\section{Exercises}
\label{ch01:exercises}

\begin{exercise}[Classify the goal]
For each scenario, name which of the four goals (confidentiality, integrity, authentication, non-repudiation) is the \emph{primary} one at stake, and name any secondary goals also required.
(a) A software vendor publishes a download alongside a value so users can check the file was not corrupted or swapped.
(b) You want to email a colleague a document so that no eavesdropper can read it.
(c) A bank needs to prove in court that a customer authorized a wire transfer.
(d) A login server must be sure the client connecting really holds the account password and is not an impostor.
\end{exercise}

\begin{exercise}[Break a shift cipher by hand]
You intercept the ciphertext \texttt{WKHUHLVQRVHFXULWB} and know only that a shift cipher (\S\ref{ch01:tiny-cipher}) was used.
(a) Recover the key $k$ and the plaintext by exhaustive search over the $25$ useful shifts. (Hint: the plaintext is an English sentence, spaces removed.)
(b) Explain why the size of the keyspace gave you no protection here.
(c) The substitution cipher has $26! \approx 2^{88}$ keys, yet is also broken easily. Name the single property both ciphers share that makes frequency analysis work, and describe what a cipher's output must look like to defeat it.
\end{exercise}

\begin{exercise}[Deterministic encryption is doomed]
Using the indistinguishability game of Figure~\ref{ch01:fig-ind-game} and Example~\ref{ch01:ex-deterministic}, prove that \emph{any} deterministic encryption scheme fails the game whenever the adversary can obtain the encryption of at least one message of its choosing. State the adversary's exact strategy and compute its winning probability. Then explain in one sentence what randomization buys you.
\end{exercise}

\begin{exercise}[Eve versus Mallory]
Give a concrete example of a system that is secure against Eve (a passive eavesdropper) but completely broken by Mallory (an active attacker). Describe exactly what Mallory does. Then explain in one sentence why confidentiality alone never implies integrity.
\end{exercise}

\begin{exercise}[Refute the obscurity argument]
A startup advertises: ``Our encryption is unbreakable because our algorithm is a closely guarded trade secret that no attacker has ever seen.'' Write a short rebuttal grounded in Kerckhoffs's principle. Give at least two distinct concrete failure modes their argument ignores, and explain why a published algorithm can be \emph{more} trustworthy than a secret one.
\end{exercise}

\begin{exercise}[Feel the numbers]
A key-search machine tries $10^{12}$ keys per second (a trillion per second, far beyond any real device per chip).
(a) How many years to exhaust a 64-bit keyspace ($2^{64}$ keys)?
(b) How many to exhaust an 80-bit keyspace?
(c) How many to exhaust a 128-bit keyspace, and how does that compare to the age of the universe ($\approx 1.4\times 10^{10}$ years)?
Use $2^{10}\approx 10^{3}$ and one year $\approx 3.15\times10^{7}$ seconds. What do your answers say about why standards moved from 64-bit to 128-bit keys?
\end{exercise}

\begin{exercise}[Break the reused keystream by hand]
Two 6-bit messages are encrypted with the \emph{same} 6-bit keystream $k$ by XOR. You intercept
$c_1 = 101101$ and $c_2 = 110011$.
(a) Compute $c_1 \xor c_2$ and explain why the keystream $k$ has disappeared from the result.
(b) Suppose you also learn that $m_1 = 011010$. Recover both $k$ and $m_2$.
(c) State, in terms of the threat model, exactly which assumption the encryptor violated.
\end{exercise}

\begin{exercise}[Rank the attack models]
Order these four attacker capabilities from weakest to strongest, and for each give a one-line real-world scenario in which an attacker would actually have that capability: ciphertext-only, known-plaintext, chosen-plaintext, chosen-ciphertext. Then explain why proving security against the \emph{strongest} model is the conservative choice.
\end{exercise}

\begin{exercise}[Why ``can't recover the plaintext'' is too weak]
Consider an encryption scheme that perfectly hides every \emph{character} of a message except that the ciphertext always leaks the message's exact length and whether the first character is a vowel. Argue that this scheme could satisfy ``the attacker cannot recover the full plaintext'' yet still be unacceptable. Tie your answer to why cryptographers use \emph{indistinguishability} (the attacker cannot tell $m_0$ from $m_1$) as the goal instead.
\end{exercise}

\begin{exercise}[Reductions, conceptually]
Explain in your own words why a theorem of the form ``if factoring is hard, then this scheme is secure'' gives more confidence than ``I tried hard to break this scheme and could not.'' In your explanation, identify what is being \emph{assumed} and what is being \emph{proven}, and state what we would learn about the world if someone exhibited an efficient attack on the scheme.
\end{exercise}

\begin{exercise}[Hybrid reasoning]
HTTPS uses slow public-key cryptography to establish a shared key, then fast symmetric cryptography (AES) for the actual data.
(a) Why not use public-key cryptography for the whole conversation?
(b) Why not use symmetric cryptography for the whole conversation, skipping public-key entirely?
(c) What problem does the public-key step solve that pure symmetric cryptography cannot, for two parties who have never met?
\end{exercise}

\begin{exercise}[Entropy versus length]
A system generates ``$256$-bit keys,'' but a flaw causes the generator to seed itself only from the current time in whole seconds, drawn from a known $30$-day window.
(a) Roughly how many distinct keys can the system actually produce? Express it as a power of two.
(b) How many bits of \emph{effective} security does the $256$-bit key really provide?
(c) An attacker can test $10^{9}$ keys per second. About how long to try them all? Relate your answer to why \S\ref{ch01:randomness} insists that key \emph{length} and key \emph{entropy} are different things.
\end{exercise}

\begin{exercise}[Order matters]
Alice wants both confidentiality and integrity and considers the three combinations in \S\ref{ch01:composition}: encrypt-and-MAC, MAC-then-encrypt, and encrypt-then-MAC.
(a) For MAC-then-encrypt, explain precisely why Bob is forced to run decryption on attacker-supplied data before he can detect tampering, and connect this to the padding-oracle attack.
(b) For encrypt-then-MAC, explain why a forged ciphertext never reaches the decryption routine.
(c) State the general principle this illustrates about composing secure building blocks, and name the modern primitive (a single interface) that removes the choice from implementers.
\end{exercise}

\begin{exercise}[Apply the mindset]
A messaging app claims end-to-end encryption. List five questions you would ask, drawn from \S\ref{ch01:mindset}, before believing the messages are actually protected against a determined adversary. For each question, briefly state what a \emph{bad} answer would reveal about a hidden weakness.
\end{exercise}

\begin{exercise}[Side channels versus the model]
A server runs a provably CPA-secure cipher, but it decrypts incoming ciphertexts and replies with either ``ok'' or ``bad padding,'' and the ``bad padding'' reply comes back measurably faster.
(a) Which attack model from \S\ref{ch01:threat-models} does this behavior hand the attacker, regardless of what the proof assumed?
(b) Explain why the security proof can be entirely correct while the deployment is still broken.
(c) Name the general defense that would close this gap.
\end{exercise}
